\IfFileExists{stacks-project.cls}{%
\documentclass{stacks-project}
}{%
\documentclass{amsart}
}

% For dealing with references we use the comment environment
\usepackage{verbatim}
\newenvironment{reference}{\comment}{\endcomment}
%\newenvironment{reference}{}{}
\newenvironment{slogan}{\comment}{\endcomment}
\newenvironment{history}{\comment}{\endcomment}

% For commutative diagrams we use Xy-pic
\usepackage[all]{xy}

% We use 2cell for 2-commutative diagrams.
\xyoption{2cell}
\UseAllTwocells

% We use multicol for the list of chapters between chapters
\usepackage{multicol}

% This is generally recommended for better output
\usepackage{lmodern}
\usepackage[T1]{fontenc}

% For cross-file-references
\usepackage{xr-hyper}

% Package for hypertext links:
\usepackage{hyperref}

% For any local file, say "hello.tex" you want to link to please
% use \externaldocument[hello-]{hello}
\externaldocument[introduction-]{introduction}
\externaldocument[conventions-]{conventions}
\externaldocument[sets-]{sets}
\externaldocument[categories-]{categories}
\externaldocument[topology-]{topology}
\externaldocument[sheaves-]{sheaves}
\externaldocument[sites-]{sites}
\externaldocument[stacks-]{stacks}
\externaldocument[fields-]{fields}
\externaldocument[algebra-]{algebra}
\externaldocument[brauer-]{brauer}
\externaldocument[homology-]{homology}
\externaldocument[derived-]{derived}
\externaldocument[simplicial-]{simplicial}
\externaldocument[more-algebra-]{more-algebra}
\externaldocument[smoothing-]{smoothing}
\externaldocument[modules-]{modules}
\externaldocument[sites-modules-]{sites-modules}
\externaldocument[injectives-]{injectives}
\externaldocument[cohomology-]{cohomology}
\externaldocument[sites-cohomology-]{sites-cohomology}
\externaldocument[dga-]{dga}
\externaldocument[dpa-]{dpa}
\externaldocument[sdga-]{sdga}
\externaldocument[hypercovering-]{hypercovering}
\externaldocument[schemes-]{schemes}
\externaldocument[constructions-]{constructions}
\externaldocument[properties-]{properties}
\externaldocument[morphisms-]{morphisms}
\externaldocument[coherent-]{coherent}
\externaldocument[divisors-]{divisors}
\externaldocument[limits-]{limits}
\externaldocument[varieties-]{varieties}
\externaldocument[topologies-]{topologies}
\externaldocument[descent-]{descent}
\externaldocument[perfect-]{perfect}
\externaldocument[more-morphisms-]{more-morphisms}
\externaldocument[flat-]{flat}
\externaldocument[groupoids-]{groupoids}
\externaldocument[more-groupoids-]{more-groupoids}
\externaldocument[etale-]{etale}
\externaldocument[chow-]{chow}
\externaldocument[intersection-]{intersection}
\externaldocument[pic-]{pic}
\externaldocument[weil-]{weil}
\externaldocument[adequate-]{adequate}
\externaldocument[dualizing-]{dualizing}
\externaldocument[duality-]{duality}
\externaldocument[discriminant-]{discriminant}
\externaldocument[derham-]{derham}
\externaldocument[local-cohomology-]{local-cohomology}
\externaldocument[algebraization-]{algebraization}
\externaldocument[curves-]{curves}
\externaldocument[resolve-]{resolve}
\externaldocument[models-]{models}
\externaldocument[functors-]{functors}
\externaldocument[equiv-]{equiv}
\externaldocument[pione-]{pione}
\externaldocument[etale-cohomology-]{etale-cohomology}
\externaldocument[proetale-]{proetale}
\externaldocument[relative-cycles-]{relative-cycles}
\externaldocument[more-etale-]{more-etale}
\externaldocument[trace-]{trace}
\externaldocument[crystalline-]{crystalline}
\externaldocument[spaces-]{spaces}
\externaldocument[spaces-properties-]{spaces-properties}
\externaldocument[spaces-morphisms-]{spaces-morphisms}
\externaldocument[decent-spaces-]{decent-spaces}
\externaldocument[spaces-cohomology-]{spaces-cohomology}
\externaldocument[spaces-limits-]{spaces-limits}
\externaldocument[spaces-divisors-]{spaces-divisors}
\externaldocument[spaces-over-fields-]{spaces-over-fields}
\externaldocument[spaces-topologies-]{spaces-topologies}
\externaldocument[spaces-descent-]{spaces-descent}
\externaldocument[spaces-perfect-]{spaces-perfect}
\externaldocument[spaces-more-morphisms-]{spaces-more-morphisms}
\externaldocument[spaces-flat-]{spaces-flat}
\externaldocument[spaces-groupoids-]{spaces-groupoids}
\externaldocument[spaces-more-groupoids-]{spaces-more-groupoids}
\externaldocument[bootstrap-]{bootstrap}
\externaldocument[spaces-pushouts-]{spaces-pushouts}
\externaldocument[spaces-chow-]{spaces-chow}
\externaldocument[groupoids-quotients-]{groupoids-quotients}
\externaldocument[spaces-more-cohomology-]{spaces-more-cohomology}
\externaldocument[spaces-simplicial-]{spaces-simplicial}
\externaldocument[spaces-duality-]{spaces-duality}
\externaldocument[formal-spaces-]{formal-spaces}
\externaldocument[restricted-]{restricted}
\externaldocument[spaces-resolve-]{spaces-resolve}
\externaldocument[formal-defos-]{formal-defos}
\externaldocument[defos-]{defos}
\externaldocument[cotangent-]{cotangent}
\externaldocument[examples-defos-]{examples-defos}
\externaldocument[algebraic-]{algebraic}
\externaldocument[examples-stacks-]{examples-stacks}
\externaldocument[stacks-sheaves-]{stacks-sheaves}
\externaldocument[criteria-]{criteria}
\externaldocument[artin-]{artin}
\externaldocument[quot-]{quot}
\externaldocument[stacks-properties-]{stacks-properties}
\externaldocument[stacks-morphisms-]{stacks-morphisms}
\externaldocument[stacks-limits-]{stacks-limits}
\externaldocument[stacks-cohomology-]{stacks-cohomology}
\externaldocument[stacks-perfect-]{stacks-perfect}
\externaldocument[stacks-introduction-]{stacks-introduction}
\externaldocument[stacks-more-morphisms-]{stacks-more-morphisms}
\externaldocument[stacks-geometry-]{stacks-geometry}
\externaldocument[moduli-]{moduli}
\externaldocument[moduli-curves-]{moduli-curves}
\externaldocument[examples-]{examples}
\externaldocument[exercises-]{exercises}
\externaldocument[guide-]{guide}
\externaldocument[desirables-]{desirables}
\externaldocument[coding-]{coding}
\externaldocument[obsolete-]{obsolete}
\externaldocument[fdl-]{fdl}
\externaldocument[index-]{index}

% Theorem environments.
%
\theoremstyle{plain}
\newtheorem{theorem}[subsection]{Theorem}
\newtheorem{proposition}[subsection]{Proposition}
\newtheorem{lemma}[subsection]{Lemma}

\theoremstyle{definition}
\newtheorem{definition}[subsection]{Definition}
\newtheorem{example}[subsection]{Example}
\newtheorem{exercise}[subsection]{Exercise}
\newtheorem{situation}[subsection]{Situation}

\theoremstyle{remark}
\newtheorem{remark}[subsection]{Remark}
\newtheorem{remarks}[subsection]{Remarks}

\numberwithin{equation}{subsection}

% Macros
%
\def\lim{\mathop{\mathrm{lim}}\nolimits}
\def\colim{\mathop{\mathrm{colim}}\nolimits}
\def\Spec{\mathop{\mathrm{Spec}}}
\def\Hom{\mathop{\mathrm{Hom}}\nolimits}
\def\Ext{\mathop{\mathrm{Ext}}\nolimits}
\def\SheafHom{\mathop{\mathcal{H}\!\mathit{om}}\nolimits}
\def\SheafExt{\mathop{\mathcal{E}\!\mathit{xt}}\nolimits}
\def\Sch{\mathit{Sch}}
\def\Mor{\mathop{\mathrm{Mor}}\nolimits}
\def\Ob{\mathop{\mathrm{Ob}}\nolimits}
\def\Sh{\mathop{\mathit{Sh}}\nolimits}
\def\NL{\mathop{N\!L}\nolimits}
\def\CH{\mathop{\mathrm{CH}}\nolimits}
\def\proetale{{pro\text{-}\acute{e}tale}}
\def\etale{{\acute{e}tale}}
\def\QCoh{\mathit{QCoh}}
\def\Ker{\mathop{\mathrm{Ker}}}
\def\Im{\mathop{\mathrm{Im}}}
\def\Coker{\mathop{\mathrm{Coker}}}
\def\Coim{\mathop{\mathrm{Coim}}}

% Boxtimes
%
\DeclareMathSymbol{\boxtimes}{\mathbin}{AMSa}{"02}

%
% Macros for moduli stacks/spaces
%
\def\QCohstack{\mathcal{QC}\!\mathit{oh}}
\def\Cohstack{\mathcal{C}\!\mathit{oh}}
\def\Spacesstack{\mathcal{S}\!\mathit{paces}}
\def\Quotfunctor{\mathrm{Quot}}
\def\Hilbfunctor{\mathrm{Hilb}}
\def\Curvesstack{\mathcal{C}\!\mathit{urves}}
\def\Polarizedstack{\mathcal{P}\!\mathit{olarized}}
\def\Complexesstack{\mathcal{C}\!\mathit{omplexes}}
% \Pic is the operator that assigns to X its picard group, usage \Pic(X)
% \Picardstack_{X/B} denotes the Picard stack of X over B
% \Picardfunctor_{X/B} denotes the Picard functor of X over B
\def\Pic{\mathop{\mathrm{Pic}}\nolimits}
\def\Picardstack{\mathcal{P}\!\mathit{ic}}
\def\Picardfunctor{\mathrm{Pic}}
\def\Deformationcategory{\mathcal{D}\!\mathit{ef}}
\usepackage{cleveref}
\crefname{originaldefinition}{original definition}{original definitions}
\Crefname{originaldefinition}{Original definition}{Original definitions}
\makeatletter
\newlabel{algebra-definition-universally-injective}{{058I}{source}{}{sourcepointer-1}{}}
\newlabel{algebra-definition-universally-injective@cref}{{[originaldefinition][0][]058I}{[1][0][]source}{}{}{}}
\makeatother
\crefname{originallemma}{original lemma}{original lemmas}
\Crefname{originallemma}{Original lemma}{Original lemmas}
\makeatletter
\newlabel{lemma-compare-etale-zariski-flat}{{0GNA}{source}{}{sourcepointer-2}{}}
\newlabel{lemma-compare-etale-zariski-flat@cref}{{[originallemma][0][]0GNA}{[1][0][]source}{}{}{}}
\makeatother
\crefname{originallemma}{original lemma}{original lemmas}
\Crefname{originallemma}{Original lemma}{Original lemmas}
\makeatletter
\newlabel{morphisms-lemma-etale-flat}{{02GS}{source}{}{sourcepointer-3}{}}
\newlabel{morphisms-lemma-etale-flat@cref}{{[originallemma][0][]02GS}{[1][0][]source}{}{}{}}
\makeatother
\crefname{originalproposition}{original proposition}{original propositions}
\Crefname{originalproposition}{Original proposition}{Original propositions}
\makeatletter
\newlabel{proposition-equivalence-quasi-coherent}{{03DX}{source}{}{sourcepointer-4}{}}
\newlabel{proposition-equivalence-quasi-coherent@cref}{{[originalproposition][0][]03DX}{[1][0][]source}{}{}{}}
\makeatother
\crefname{originalproposition}{original proposition}{original propositions}
\Crefname{originalproposition}{Original proposition}{Original propositions}
\makeatletter
\newlabel{proposition-same-cohomology-quasi-coherent}{{03DW}{source}{}{sourcepointer-5}{}}
\newlabel{proposition-same-cohomology-quasi-coherent@cref}{{[originalproposition][0][]03DW}{[1][0][]source}{}{}{}}
\makeatother
\crefname{originalremark}{original remark}{original remarks}
\Crefname{originalremark}{Original remark}{Original remarks}
\makeatletter
\newlabel{remark-change-topologies-ringed-sites}{{03FH}{source}{}{sourcepointer-6}{}}
\newlabel{remark-change-topologies-ringed-sites@cref}{{[originalremark][0][]03FH}{[1][0][]source}{}{}{}}
\makeatother
\crefname{originalsection}{original section}{original sections}
\Crefname{originalsection}{Original section}{Original sections}
\makeatletter
\newlabel{schemes-section-quasi-coherent}{{01LA}{source}{}{sourcepointer-7}{}}
\newlabel{schemes-section-quasi-coherent@cref}{{[originalsection][0][]01LA}{[1][0][]source}{}{}{}}
\makeatother
\crefname{originallemma}{original lemma}{original lemmas}
\Crefname{originallemma}{Original lemma}{Original lemmas}
\makeatletter
\newlabel{sites-modules-lemma-exactness-pushforward-pullback}{{03DC}{source}{}{sourcepointer-8}{}}
\newlabel{sites-modules-lemma-exactness-pushforward-pullback@cref}{{[originallemma][0][]03DC}{[1][0][]source}{}{}{}}
\makeatother
\crefname{originallemma}{original lemma}{original lemmas}
\Crefname{originallemma}{Original lemma}{Original lemmas}
\makeatletter
\newlabel{sites-modules-lemma-pullback-internal-hom}{{0GN2}{source}{}{sourcepointer-9}{}}
\newlabel{sites-modules-lemma-pullback-internal-hom@cref}{{[originallemma][0][]0GN2}{[1][0][]source}{}{}{}}
\makeatother
\crefname{originallemma}{original lemma}{original lemmas}
\Crefname{originallemma}{Original lemma}{Original lemmas}
\makeatletter
\newlabel{sites-modules-lemma-tensor-product-pullback}{{03EL}{source}{}{sourcepointer-10}{}}
\newlabel{sites-modules-lemma-tensor-product-pullback@cref}{{[originallemma][0][]03EL}{[1][0][]source}{}{}{}}
\makeatother
\title{Properties of the functor from quasi-coherent modules to modules on a big or small site}
\author{The Stacks Project Authors}
\date{stacks-project:descent.tex:local-snapshot-20261012; GNU FDL 1.2 or later}
\begin{document}
\maketitle
\noindent\textit{Editorial note.} All delivered mathematical text is the author's own, copied from the pinned Stacks Project chapter file descent.tex; no statement or proof step was rewritten, reordered or removed.\par\medskip
\noindent\textit{Editorial note.} The statement carries the author's own footnote warning that part (10) is misleading in view of part (6); the footnote is inside the counted statement span and is reproduced verbatim.\par\medskip
\noindent\textit{Editorial note.} The counted spans carry 10 navigation stubs, one per cited result; none of them replaces a load-bearing step, since the counterexamples and the final reduction are written out in the span. No bibliography citation occurs in the spans.\par\medskip
\noindent\textit{Editorial note.} The preamble span points at the chapter's own preamble.tex shipped with the frozen source, used verbatim as staged.\par\medskip
\noindent\textit{Allowed standard prerequisites (editorial).}\begin{itemize}
\item The functor sending a quasi-coherent module to the module it represents on a tau-site. Same chapter, Proposition proposition-equivalence-quasi-coherent, invoked at proof lines 3310-3311 and named in the statement. Supplies the functors whose properties are listed and part (1).
\item The construction as a pullback of ringed sites and the exactness properties of such pullbacks. Same chapter, Remark remark-change-topologies-ringed-sites and Modules on Sites, Lemma sites-modules-lemma-exactness-pushforward-pullback, invoked at proof lines 3317-3322. Supplies commutation with colimits and hence right exactness.
\item Etale morphisms are flat, and pullback along flat morphisms of ringed sites commutes with internal hom out of a finitely presented module. Morphisms, Lemma morphisms-lemma-etale-flat and Modules on Sites, Lemmas sites-modules-lemma-tensor-product-pullback and sites-modules-lemma-pullback-internal-hom, invoked at proof lines 3330-3362. Supplies parts (5), (7) and (8).
\item Cohomology of a quasi-coherent module agrees on the small and big sites on affines. Same chapter, Proposition proposition-same-cohomology-quasi-coherent, invoked at proof lines 3385-3388. Supplies the vanishing of the first cohomology used in part (10).
\end{itemize}\medskip
\ifdefined\hypertarget\hypertarget{sourcepointer-1}{}\fi
\noindent\textit{Source reference.} Original definition 058I: sources/stacks-chow/tags.txt, source lines 6781--6781. Pointer outside the counted span, invoked once for context; the load-bearing argument of the counted statement is the author's proof reproduced in full in the proof span.\par\medskip
\ifdefined\hypertarget\hypertarget{sourcepointer-2}{}\fi
\noindent\textit{Source reference.} Original lemma 0GNA: sources/stacks-descent/pinned-original-source.tex, source lines 3221--3221. Pointer outside the counted span, invoked once for context; the load-bearing argument of the counted statement is the author's proof reproduced in full in the proof span.\par\medskip
\ifdefined\hypertarget\hypertarget{sourcepointer-3}{}\fi
\noindent\textit{Source reference.} Original lemma 02GS: sources/stacks-chow/tags.txt, source lines 3239--3239. Pointer outside the counted span, invoked once for context; the load-bearing argument of the counted statement is the author's proof reproduced in full in the proof span.\par\medskip
\ifdefined\hypertarget\hypertarget{sourcepointer-4}{}\fi
\noindent\textit{Source reference.} Original proposition 03DX: sources/stacks-descent/pinned-original-source.tex, source lines 2716--2716. Pointer outside the counted span, invoked once for context; the load-bearing argument of the counted statement is the author's proof reproduced in full in the proof span.\par\medskip
\ifdefined\hypertarget\hypertarget{sourcepointer-5}{}\fi
\noindent\textit{Source reference.} Original proposition 03DW: sources/stacks-descent/pinned-original-source.tex, source lines 2990--2990. Pointer outside the counted span, invoked once for context; the load-bearing argument of the counted statement is the author's proof reproduced in full in the proof span.\par\medskip
\ifdefined\hypertarget\hypertarget{sourcepointer-6}{}\fi
\noindent\textit{Source reference.} Original remark 03FH: sources/stacks-descent/pinned-original-source.tex, source lines 2587--2587. Pointer outside the counted span, invoked once for context; the load-bearing argument of the counted statement is the author's proof reproduced in full in the proof span.\par\medskip
\ifdefined\hypertarget\hypertarget{sourcepointer-7}{}\fi
\noindent\textit{Source reference.} Original section 01LA: sources/stacks-chow/tags.txt, source lines 2079--2079. Pointer outside the counted span, invoked once for context; the load-bearing argument of the counted statement is the author's proof reproduced in full in the proof span.\par\medskip
\ifdefined\hypertarget\hypertarget{sourcepointer-8}{}\fi
\noindent\textit{Source reference.} Original lemma 03DC: sources/stacks-chow/tags.txt, source lines 4415--4415. Pointer outside the counted span, invoked once for context; the load-bearing argument of the counted statement is the author's proof reproduced in full in the proof span.\par\medskip
\ifdefined\hypertarget\hypertarget{sourcepointer-9}{}\fi
\noindent\textit{Source reference.} Original lemma 0GN2: sources/stacks-chow/tags.txt, source lines 20777--20777. Pointer outside the counted span, invoked once for context; the load-bearing argument of the counted statement is the author's proof reproduced in full in the proof span.\par\medskip
\ifdefined\hypertarget\hypertarget{sourcepointer-10}{}\fi
\noindent\textit{Source reference.} Original lemma 03EL: sources/stacks-chow/tags.txt, source lines 4460--4460. Pointer outside the counted span, invoked once for context; the load-bearing argument of the counted statement is the author's proof reproduced in full in the proof span.\par\medskip
\begin{lemma}
\label{lemma-equivalence-quasi-coherent-limits}
Let $S$ be a scheme.
Let $\tau \in \{Zariski, \linebreak[0] \etale,
\linebreak[0] smooth, \linebreak[0] syntomic, \linebreak[0] fppf\}$.
The functors
$$
\QCoh(\mathcal{O}_S)
\longrightarrow
\textit{Mod}((\Sch/S)_\tau, \mathcal{O})
\quad\text{and}\quad
\QCoh(\mathcal{O}_S)
\longrightarrow
\textit{Mod}(S_\tau, \mathcal{O})
$$
defined by the rule $\mathcal{F} \mapsto \mathcal{F}^a$ seen in
Proposition \ref{proposition-equivalence-quasi-coherent}
are
\begin{enumerate}
\item fully faithful,
\item commmute with direct sums,
\item commmute with colimits,
\item right exact,
\item exact as a functor
$\QCoh(\mathcal{O}_S) \to \textit{Mod}(S_\etale, \mathcal{O})$,
\item {\bf not} exact as a functor
$\QCoh(\mathcal{O}_S) \to
\textit{Mod}((\Sch/S)_\tau, \mathcal{O})$
in general,
\item given two quasi-coherent $\mathcal{O}_S$-modules
$\mathcal{F}$, $\mathcal{G}$ we have
$(\mathcal{F} \otimes_{\mathcal{O}_S} \mathcal{G})^a =
\mathcal{F}^a \otimes_\mathcal{O} \mathcal{G}^a$,
\item if $\tau = \etale$ or $\tau = Zariski$,
given two quasi-coherent $\mathcal{O}_S$-modules
$\mathcal{F}$, $\mathcal{G}$ such that $\mathcal{F}$
is of finite presentation we have
$(\SheafHom_{\mathcal{O}_S}(\mathcal{F}, \mathcal{G}))^a =
\SheafHom_\mathcal{O}(\mathcal{F}^a, \mathcal{G}^a)$ in
$\textit{Mod}(S_\tau, \mathcal{O})$,
\item given two quasi-coherent $\mathcal{O}_S$-modules
$\mathcal{F}$, $\mathcal{G}$ we do {\bf not} have
$(\SheafHom_{\mathcal{O}_S}(\mathcal{F}, \mathcal{G}))^a =
\SheafHom_\mathcal{O}(\mathcal{F}^a, \mathcal{G}^a)$
in $\textit{Mod}((\Sch/S)_\tau, \mathcal{O})$ in general
even if $\mathcal{F}$ is of finite presentation, and
\item given a short exact sequence
$0 \to \mathcal{F}_1^a \to \mathcal{E} \to \mathcal{F}_2^a \to 0$
of $\mathcal{O}$-modules then $\mathcal{E}$ is
quasi-coherent\footnote{Warning: This is misleading. See part (6).}, i.e.,
$\mathcal{E}$ is in the essential image of the functor.
\end{enumerate}
\end{lemma}

\begin{proof}
Part (1) we saw in
Proposition \ref{proposition-equivalence-quasi-coherent}.

\medskip\noindent
We have seen in
Schemes, Section \ref{schemes-section-quasi-coherent}
that a colimit of quasi-coherent sheaves on a scheme is a quasi-coherent
sheaf. Moreover, in
Remark \ref{remark-change-topologies-ringed-sites}
we saw that $\mathcal{F} \mapsto \mathcal{F}^a$ is the pullback functor
for a morphism of ringed sites, hence commutes with all colimits, see
Modules on Sites, Lemma
\ref{sites-modules-lemma-exactness-pushforward-pullback}.
Thus (3) and its special case (2) hold.

\medskip\noindent
This also shows that the functor is right exact (i.e., commutes with
finite colimits), hence (4).

\medskip\noindent
The functor $\QCoh(\mathcal{O}_S) \to
\textit{Mod}(S_\etale, \mathcal{O})$,
$\mathcal{F} \mapsto \mathcal{F}^a$
is left exact because an \'etale morphism is flat, see
Morphisms, Lemma \ref{morphisms-lemma-etale-flat}.
This proves (5).

\medskip\noindent
To see (6), suppose that $S = \Spec(\mathbf{Z})$.
Then $2 : \mathcal{O}_S \to \mathcal{O}_S$ is injective but the associated
map of $\mathcal{O}$-modules on $(\Sch/S)_\tau$ isn't
injective because $2 : \mathbf{F}_2 \to \mathbf{F}_2$ isn't injective
and $\Spec(\mathbf{F}_2)$ is an object of $(\Sch/S)_\tau$.

\medskip\noindent
Part (7) holds because, as mentioned above, the functor
$\mathcal{F} \mapsto \mathcal{F}^a$ is the pullback functor
for a morphism of ringed sites and such commute with tensor
products by Modules on Sites, Lemma
\ref{sites-modules-lemma-tensor-product-pullback}.

\medskip\noindent
Part (8) is obvious if $\tau = Zariski$ because the category of
$\mathcal{O}$-modules on $S_{Zar}$ is the same as the category
of $\mathcal{O}_S$-modules on the topological space $S$.
If $\tau = \etale$ then (8) holds because, as mentioned above,
the functor $\mathcal{F} \mapsto \mathcal{F}^a$ is the pullback functor
for the flat morphism of ringed sites
$(S_\etale, \mathcal{O}) \to (S_{Zar}, \mathcal{O}_S)$, see
Lemma \ref{lemma-compare-etale-zariski-flat}.
Pullback by flat morphisms of ringed sites commutes with
taking internal hom out of a finitely presented module by
Modules on Sites, Lemma \ref{sites-modules-lemma-pullback-internal-hom}.

\medskip\noindent
To see (9), suppose that $S = \Spec(\mathbf{Z})$. Let
$\mathcal{F} = \Coker(2 : \mathcal{O}_S \to \mathcal{O}_S)$
and $\mathcal{G} = \mathcal{O}_S$.
Then $\mathcal{F}^a = \Coker(2 : \mathcal{O} \to \mathcal{O})$
and $\mathcal{G}^a = \mathcal{O}$.
Hence
$\SheafHom_\mathcal{O}(\mathcal{F}^a, \mathcal{G}^a) = \mathcal{O}[2]$
is equal to the $2$-torsion in $\mathcal{O}$, which is not zero,
see proof of (6). On the other hand, the module
$\SheafHom_{\mathcal{O}_S}(\mathcal{F}, \mathcal{G})$
is zero.

\medskip\noindent
Proof of (10).
Let $0 \to \mathcal{F}_1^a \to \mathcal{E} \to \mathcal{F}_2^a \to 0$
be a short exact sequence of $\mathcal{O}$-modules with $\mathcal{F}_1$
and $\mathcal{F}_2$ quasi-coherent on $S$. Consider the restriction
$$
0 \to \mathcal{F}_1 \to \mathcal{E}|_{S_{Zar}} \to \mathcal{F}_2
$$
to $S_{Zar}$. By
Proposition \ref{proposition-same-cohomology-quasi-coherent}
we see that on any affine $U \subset S$ we have
$H^1(U, \mathcal{F}_1^a) = H^1(U, \mathcal{F}_1) = 0$.
Hence the sequence above is also exact on the right. By
Schemes, Section \ref{schemes-section-quasi-coherent}
we conclude that $\mathcal{F} = \mathcal{E}|_{S_{Zar}}$ is
quasi-coherent. Thus we obtain a commutative diagram
$$
\xymatrix{
& \mathcal{F}_1^a \ar[r] \ar[d] &
\mathcal{F}^a \ar[r] \ar[d] &
\mathcal{F}_2^a \ar[r] \ar[d] & 0 \\
0 \ar[r] &
\mathcal{F}_1^a \ar[r] &
\mathcal{E} \ar[r] &
\mathcal{F}_2^a \ar[r] & 0
}
$$
To finish the proof it suffices to show that the top row is also
right exact. To do this, denote once more $U = \Spec(A) \subset S$
an affine open of $S$. We have seen above that
$0 \to \mathcal{F}_1(U) \to \mathcal{E}(U) \to \mathcal{F}_2(U) \to 0$
is exact. For any affine scheme $V/U$,
$V = \Spec(B)$ the map $\mathcal{F}_1^a(V) \to \mathcal{E}(V)$
is injective. We have $\mathcal{F}_1^a(V) = \mathcal{F}_1(U) \otimes_A B$
by definition. The injection
$\mathcal{F}_1^a(V) \to \mathcal{E}(V)$ factors as
$$
\mathcal{F}_1(U) \otimes_A B \to
\mathcal{E}(U) \otimes_A B \to \mathcal{E}(V)
$$
Considering $A$-algebras $B$ of the form $B = A \oplus M$
we see that $\mathcal{F}_1(U) \to \mathcal{E}(U)$ is
universally injective (see
Algebra, Definition \ref{algebra-definition-universally-injective}).
Since $\mathcal{E}(U) = \mathcal{F}(U)$ we conclude that
$\mathcal{F}_1 \to \mathcal{F}$ remains injective after any base change,
or equivalently that $\mathcal{F}_1^a \to \mathcal{F}^a$ is injective.
\end{proof}

\end{document}
